\section[Recursive numerical ranges]{A spectral edge floor and simultaneous recursive numerical ranges}
\label{sec:op3-recursive-vwf-ranges}

\paragraph{Scope.}
This section supplies an implemented preconditioner modification and
range bounds for the mixed recursion of
\cref{cor:op3-mixed-recursive-budgets}. All graphs here are supplied,
connected, loopless positive-weight graphs; parallel edge records are
allowed. Work is exact-real word work. The new statements are
\statusproved\ drafts awaiting independent review. They do not establish
a local graph-discovery algorithm or a fast preconditioner constructor.

\begin{lemma}[A paid spectral preconditioner edge floor; \statusproved]
\label{lem:op3-spectral-preconditioner-floor}
Let $G\preceq Q\preceq\kappa_0G$ be graph Laplacians on $n\geq2$
vertices. Let $c$ be the smallest positive edge weight of $G$, and let
$m_Q$ be the number of edge records of $Q$. Delete exactly the edges of
$Q$ whose weights are strictly below
\[
 t=\frac{c}{2m_Q(n-1)},
\]
and double every retained weight, obtaining $H$. Then
\[
 G\preceq H\preceq2\kappa_0G,
 \qquad c_H\geq\frac{c}{m_Q(n-1)},
 \qquad W_H\leq2\kappa_0W_G,
\]
where $W$ denotes the sum of edge weights. The complete modification
costs $O(n+m_G+m_Q)$ work and $O(n+m_Q)$ newly allocated words,
including retained and deleted records.
If $Q$ is a supplied forest with a retained-root core, no forest edge is
deleted and the core stays connected. The original objective is unchanged.
\end{lemma}
\begin{proof}
A simple path in $G$ and weighted Cauchy--Schwarz give
$(x_u-x_v)^2\leq(n-1)x^TGx/c$ for any pair $u,v$.
The deleted Laplacian $P$ has total weight at most $m_Qt$, so
$P\preceq G/2$. Therefore $H=2(Q-P)\succeq G$ and
$H\preceq2\kappa_0G$. The edge floor is the retained threshold after
doubling. Taking traces proves the weight-sum bound, since a graph
Laplacian has trace twice its total edge weight.

Every forest edge is a bridge of $Q$. For its component cut $S$, the
indicator-vector inequality $Q\succeq G$ implies that its weight is at
least $G(S,V\setminus S)\geq c$. Thus it is not deleted. Connectivity
of $H$ follows from $H\succeq G$, and with all forest edges retained
this also gives connectivity of its root core. Scan every original edge
for $c$, then every supplied edge for the threshold comparison, copying
and scaling retained records and charging deleted records. No spectral
test is part of the modification: the input order is its contract.
The sparse matrix $G$ and its vertex functions are untouched, so there
is no objective-pruning error to budget.
\end{proof}

For a VWF problem $\Phi$, use the following nonnegative quantities:
\[
 B\geq-\Phi^*,\quad
 S=\sum_i\max(T_i,0),\quad C=\sum_i h_i,\quad
 R=\max\{1,U,\max_i(-L_i)\},
\]
where $T_i$ is the final derivative, $h_i$ its largest curvature, and
$U$ the largest nonnegative split (zero if absent). Assume $\Phi(0)\leq0$
and finite minimum, so $\sum_iT_i\geq0$. The symbols $S,C,R$ in this
section are numerical ranges, not support sets or input piece counts.

\begin{lemma}[Range propagation for a single recursive step; \statusproved]
\label{lem:op3-single-recursive-ranges}
Let a mixed accelerated call use adjusted quality $\kappa\geq1$ and
requested additive error $\eta>0$. Its normalized inner objective $E$
satisfies
\begin{align*}
 -E^*&\leq B_I:=33\kappa^2(B+\eta),\\
 S_I&\leq S+32\kappa\sqrt{\kappa W_G(B+\eta)}.
\end{align*}
Its domains, splits and curvatures equal those of $\Phi$.
Exact forest elimination on $H$ gives a coarse objective $F$ with
\[
 -F^*\leq B_I,\quad S_F\leq S_I,\quad C_F\leq C+W_H,
 \quad U_F\leq U+nS_I/c_H.
\]
Its lower endpoints are a subset of the original endpoints.
For any canonical coarse refinement state $a$ with
$F(a)\leq F(0)\leq0$, the exact residual $F_a$ has
\begin{align*}
 -F_a^*&\leq B_I,&
 S_a&\leq S_F+4\sqrt{2W_FB_I},& C_a&=C_F,\\
 R_a&\leq2R_F+\sqrt{8K_FB_I},&&K_F=(n_F-1)/c_F.
\end{align*}
Compression with common cutoff $\tau\leq1$ and total error $\Xi$
then gives a child objective satisfying
\[
 B_{\rm child}\leq2B_I+2\Xi,\quad
 S_{\rm child}=S_a,\quad C_{\rm child}\leq C_a,
 \quad R_{\rm child}\leq2R_a.
\]
These bounds apply to every refinement call without summing its state
ranges over the number of earlier calls.
\end{lemma}
\begin{proof}
The mixed accelerated theorem gives
$\|x\|_H^2\leq64\kappa^2(\Delta+\eta)$, where
$\Delta\leq B$. Since $\Phi_x\geq\Phi$,
\[
 -E^*=\Phi_x(0)-\Phi_x^*
 \leq\tfrac12\|x\|_H^2+\Delta\leq B_I.
\]
The terminal-mass proof of
\cref{cor:op3-generic-proximal-terminal-mass} applies with
$\Delta+\eta\leq B+\eta$. Normalization adds only a linear term and
subtracts each zero-anchor value. The forest assertions follow from
\cref{prop:op3-vwf-forest-ranges}; positive parts of disjoint subtree
tail sums cannot increase their total, and the minimum energy is preserved.

Both the gap at $a$ and the gap at zero are at most $B_I$. The
gap-to-distance inequality and triangle inequality give
$\|a\|_{G_F}\leq2\sqrt{2B_I}$. Thus
\[
 \|G_Fa\|_1\leq2\sqrt{W_F}\|a\|_{G_F}
 \leq4\sqrt{2W_FB_I}.
\]
The residual final slopes add $(G_Fa)_i$, proving the mass bound.
Canonicalization gives $\min_iL_i\leq\min_i a_i\leq U_F$;
the path inequality gives diameter at most $\sqrt{8K_FB_I}$.
Hence $L_i-a_i\geq-2R_F-\sqrt{8K_FB_I}$, while each positive
shifted split is at most $U_F-\min_i a_i\leq2R_F$.
Also $F_a^*=F^*-F(a)\geq F^*$.
The residual is constructed from the original coarse functions at the
current state; previous linear shifts are not added again.

The global compression inequality
$\widehat F_a(x)\geq2F_a(x/2)-2\Xi$ bounds its minimum below by
$2F_a^*-2\Xi$. The compressor preserves final slopes and does not
increase total curvature. Every moved split has magnitude at most
$2\max(R_a,\tau)\leq2R_a$. Lower endpoints are unchanged. These
facts prove all child bounds, including zero gap and negative individual
terminal slopes. If the coarse graph has one vertex, minimize that
scalar function directly and terminate instead of defining $c_F$.
\end{proof}

\begin{theorem}[A simultaneous logarithmic range bound; \statusproved]
\label{thm:op3-simultaneous-recursive-ranges}
Consider any depth-$D$ execution satisfying the preceding contracts,
using the spectral floor before every exact forest pass and
\cref{cor:op3-mixed-recursive-budgets}'s error schedule.
Let $P\geq2$ bound all vertex counts, supplied preconditioner edge counts
and adjusted qualities. Write $Z\geq2$ for an upper bound on
\[
 R_0,\ S_0,\ C_0,\ W_0,\ c_0^{-1},\ B_0+\max(1,\eta_0),
 \ \eta_0^{-1}.
\]
Use a common compression cutoff
\[
 \tau=\min\{1,\Xi_0/C_a\},\qquad
 \Xi_0=\frac{\eta}{2^{35}\kappa},
\]
with $\tau=1$ if $C_a=0$. All quantities
$B,S,C,R,W,c^{-1},\eta^{-1},\tau^{-1}$ at any level have logarithm
at most
\[
 O\bigl((D+1)\log P+\log Z\bigr).
\]
The number of nonzero dyadic compression bins per vertex is bounded by
the same expression. This conclusion needs no lower bound on nonzero
curvature drops, gaps or split separations.
\end{theorem}
\begin{proof}
Set $Q=2^{40}P^4$ and $E_0=\max(1,\eta_0)$. Along a path, the
graph-floor and trace bounds give $c_{d+1}^{-1}\leq P^2c_d^{-1}$
and $W_{d+1}\leq PW_d$. The schedule gives
$\eta_{d+1}^{-1}\leq2^{34}P\eta_d^{-1}$ and $\eta_d\leq E_0$.
Since $\Xi\leq\Xi_0\leq\eta_d$,
\[
 B_{d+1}+E_0\leq(66P^2+3)(B_d+E_0).
\]
Consequently $B_d+E_0,W_d,c_d^{-1},\eta_d^{-1}\leq Q^dZ$.
The two terminal-mass additions in the preceding lemma give
\[
 S_{d+1}\leq S_d+65P^2\sqrt{W_d(B_d+E_0)},\qquad
 C_{d+1}\leq C_d+PW_d.
\]
Geometric summation implies $S_d,C_d\leq Q^{d+1}Z$.
The inner mass obeys $S_I\leq2Q^{d+1}Z$, and the preconditioner floor
obeys $c_H^{-1}\leq Q^{d+1}Z$. Combining the forest, residual and
compression radius bounds therefore gives inductively
\[
 R_d\leq Q^{2d+2}Z^2.
\]
For completeness, in its induction step the term $4R_d$ contributes
at most $4Q^{2d+2}Z^2$, the forest addition contributes at most
$8P Q^{2d+2}Z^2$, and the diameter term contributes at most
$33P^2Q^{2d+2}Z^2$. Their sum is below $Q^{2d+4}Z^2$.
Also $C_a\leq2Q^{d+1}Z$, whence
\[
 \tau^{-1}\leq\max\{1,2^{35}PC_a/\eta_d\}
 \leq Q^{2d+2}Z^2.
\]
Taking logarithms proves the stated bounds and the bin count
$O(1+\log(R_a/\tau))$. The chosen cutoff gives
$\Xi\leq C_a\tau^2/8\leq\Xi_0$.
The ordering of these inductions matters: graph weights and energy are
bounded first, slopes and curvature second, and radius last. No new
$A^2$ input-gap bound is substituted at each recursive level.
\end{proof}

\paragraph{Which numbers the result controls.}
The derivative-at-zero mass obeys
$\sum_i|f_i'(0)|\leq2S+CR$, because $\sum_iT_i\geq0$ and the
derivative drop from the final ray to zero is at most $h_iR$.
Individual zero-anchor values remain in $[-B,0]$ after an exact pass;
residual anchors are zero, and compression's negative anchor shifts
total at most $\Xi$. These facts also bound ordinary piece coefficients
by fixed-degree polynomials in $B,S,C,R$. Tiny nonzero jumps and affine
tag cancellations can remain. The theorem is not Assumption 3.15 for
every encountered nonzero, nor is it a rational bit-complexity theorem.
The implemented persistent and compression operations charge their
comparisons, arithmetic, stored versions and scans without depending on
the magnitude of a nonzero jump.

\paragraph{Source construction still requires a separate check.}
CPW Lemma 4.4, PDF p.~19, explicitly assumes polynomial edge-weight
ratio; its Algorithm 3 and work recurrence are on p.~20. The preceding
range theorem alone does not imply that ratio is polynomial in the
current, shrinking vertex count. It cannot silently discharge the source
hypothesis. Sections 5.1--5.6 reduce construction to weighted low-stretch
trees, exact tree decomposition/routing and core spectral sparsification.
The weighted source primitives are now reconciled in
\cref{sec:op3-bounded-core-sparsifier}, using a corrected rooted-piece
construction, and the capped supplied recurrence is completed in
\cref{thm:op3-complete-supplied-recursion}. It still needs a separate
local-discovery and cumulative-work argument for original OP3.
General OP3 remains \statusopen.
